\section{Interpretation and Bias--Variance Considerations}\label{sec:analysis}

This section does not provide a new convergence theorem. Instead, it clarifies how intrinsic token weighting should be interpreted relative to standard advantage estimation and why it can still be useful despite introducing bias.

\subsection{Credit Redistribution, Not Reward Redefinition}

The sequence-level reward remains unchanged throughout our method. What changes is the allocation of gradient mass over score-function terms:
\begin{equation}
 \nabla_\theta \log \pi_\theta(y \mid x) = \sum_{t=1}^{T} \nabla_\theta \log \pi_\theta(y_t \mid y_{<t}, x).
\end{equation}
Uniform token updates treat each summand identically once a sequence-level advantage is fixed. Intrinsic token weighting instead introduces a prior over token salience. This is best viewed as \emph{credit redistribution}: the same verifier reward is being pushed through the trajectory, but some actions are declared more update-worthy than others.

\subsection{Relationship to Advantage Estimation}

A learned critic can be interpreted as a mechanism for constructing token- or prefix-specific credit. In actor--critic methods, the update at time $t$ depends on an estimate of how much the observed trajectory outcome differs from the expected future return under the current prefix. This is powerful when the state abstraction is reliable. In language generation, however, the prefix is often a noisy and semantically unstable summary of downstream return. Intrinsic weighting offers a different inductive bias: rather than predicting expected reward from a prefix, it uses local statistics of the policy distribution itself to decide where to place more of the update.

This comparison is deliberately modest. Intrinsic weighting should not be read as a replacement for all forms of advantage estimation. Rather, it is a useful comparison point when one wants to avoid critic training, token-level reward models, or search trees while still moving beyond uniform token treatment.

\subsection{Bias Is the Price of Structured Allocation}

Because $w_t(x,y)$ depends on the sampled trajectory and the policy's own predictive distribution, the estimator
\begin{equation}
 g_w(x,y) = A(x,y) \sum_{t=1}^{T} w_t(x,y) \, \nabla_\theta \log \pi_\theta(y_t \mid y_{<t}, x)
\end{equation}
is generally biased relative to the standard REINFORCE estimator. This bias is not a bug in the analysis; it is the intended effect of introducing a structured credit prior. The practical question is therefore whether the bias is \emph{useful}: does the weighting shift optimization toward more semantically consequential decisions and away from filler tokens?

Our controlled experiments suggest that the answer can be yes, but not for the simplistic reason that weighting uniformly reduces gradient variance. In fact, the stronger weighting rules often increase our weighted within-trajectory score-variance proxy. The gains appear to come from \emph{selective concentration} rather than indiscriminate smoothing. This matters for interpretation: intrinsic token weighting is better understood as an inductive bias over update locations than as a classic unbiased variance-reduction trick.

\subsection{Why Surprisal and Entropy Reduction?}

The two weighting rules we study encode different hypotheses about token importance.

\paragraph{Surprisal.}
Low-probability sampled actions often correspond to less routine policy decisions. Weighting by surprisal therefore asks the optimizer to pay disproportionate attention to actions that the current policy considered relatively unlikely. In our arithmetic benchmark, this produces the most consistent improvement in pass@4, suggesting that rare-but-decisive symbolic choices are a useful target for extra credit.

\paragraph{Entropy reduction.}
A large entropy drop indicates that, after the current token, the policy's continuation distribution becomes more concentrated. Such tokens are plausible proxies for commitment points or branch-closing decisions. In our program-trace benchmark this signal is particularly effective, strongly increasing both success rate and the fraction of token mass assigned to reward-relevant positions.

\subsection{Failure Modes}

The same interpretation also reveals failure modes. Entropy reduction can overemphasize formatting or syntactic closure in natural-language tasks, and surprisal can overweight noisy or spurious low-probability actions. The value of intrinsic weighting is therefore task-dependent. This is one reason we prefer a controlled benchmark first: it lets us see whether the weighting actually moves mass toward known important decisions before scaling the method to larger RLVR systems.

\subsection{Takeaway}

The main conceptual takeaway is simple. Sequence-level RL for language always contains an implicit token-credit rule. Uniform weighting is one such rule, not a law of nature. Intrinsic token weighting offers another: a small, biased, but computationally cheap mechanism for redistributing update mass using policy-internal signals. Its promise is not that it solves credit assignment in general, but that it may provide a useful and modular prior in settings where critics, token rewards, or rollout trees are undesirable.
